## TEMP holding note -- does Q1's degradation require real environmental signal? (2026-08-17)

**Which RQ this closes a gap in, and why**: RQ2a item 1 (environmental conditioning trades
accuracy between subgroups: it helps the plots the shared Chapman-Richards curve fits worst and
hurts the plots it already fits best). The results draft flagged this as genuinely open: the
five-seed reseed check already run confirmed Q4's correction is reproducible across seeds
(`TEMP_rq2a_residual_reduction_results_2026-08-11.tex`), but could not tell whether Q1's
degradation reflects a real cost of using environmental information, or is simply what any
flexible model with output variance would show on Q1 regardless of what it is fed, since Q1 rows
were selected specifically because the CR baseline's own error was already smallest there.

**Method**: `models/spatial_attribution/rq2a_permutation_check.py`. Reuses
`rq2a_xgb_check.py`'s own `fit_one_arm()` and `run_reduction_check()`, and
`rq2_residual_reduction.py`'s `compute_residual_reduction()`, unmodified -- no reimplementation of
fitting or quartile logic, only a new permutation step before fitting. Set3
(`nested_set3_gated_terrain_wind_vif`), fixed XGBoost config (n_estimators=500, max_depth=4,
learning_rate=0.04), same 5-fold spatial_block_kfold split as the real env-conditioned model, both
cohorts. Before fitting each fold, the whole environmental-column block is permuted with ONE
shared row-permutation (not one shuffle per column independently, which would additionally
scramble the environmental variables' own real correlation structure) -- this keeps a real,
internally-consistent-looking environmental vector attached to every row, just the wrong plot's
vector. The model still trains on real heights and a real no-environment feature set; only the
environmental values are fake. Permutation seed = 42 + fold number, so each fold gets its own but
reproducible fake assignment.

### Results

| Cohort | Arm | Mean reduction (m) | Q1 (best-fit) | Q2 | Q3 | Q4 (worst-fit) |
|---|---|---:|---:|---:|---:|---:|
| 4survey | Real environment (`xgb_env_terrain_fixed_...`, already fit by `rq2a_xgb_check.py`) | 1.597±0.330 | -1.207 | 0.507 | 2.255 | 4.833 |
| 4survey | **Permuted environment (this check)** | 1.219±0.224 | -1.304 | 0.272 | 1.936 | 3.973 |
| 6survey | Real environment | 0.361±0.183 | -0.723 | -0.020 | 0.572 | 1.614 |
| 6survey | **Permuted environment (this check)** | 0.248±0.146 | -0.565 | -0.010 | 0.386 | 1.180 |

**Q1's degradation does not require real environmental signal.** Under permuted (scrambled,
meaningless) environmental features, Q1 still degrades by almost exactly the same amount as under
the real ones -- 4survey: -1.304m permuted vs. -1.207m real (permuted is if anything slightly
worse); 6survey: -0.565m permuted vs. -0.723m real (permuted is somewhat less bad, but the same
direction and a comparable magnitude). A model fed deliberately fake environmental values degrades
the best-fitting quartile just as much as a model fed the real ones. This is the opposite of what
"conditioning on real environment causes a real cost on Q1" would predict -- if that were the
mechanism, permuting the environment away should have made Q1's degradation shrink toward zero,
and it does not.

**Q4's correction DOES depend on real environmental signal, partially.** Permuting the
environment away shrinks Q4's correction meaningfully: 4survey 3.973m (permuted) vs. 4.833m (real),
an 18% reduction; 6survey 1.180m vs. 1.614m, a 27% reduction. The aggregate mean reduction follows
the same pattern (1.219m vs. 1.597m on 4survey, a 24% reduction; 0.248m vs. 0.361m on 6survey, a
31% reduction). Real environmental information contributes a genuine, non-trivial share of Q4's
correction -- permuting it away does not eliminate the correction entirely (a permuted-feature
model with pure output flexibility still recovers 76-82% of the real model's Q4 correction), so
Q4's benefit is not purely an artefact of model flexibility either, but real signal clearly adds
something on top of what flexibility alone provides.

**What this changes about RQ2a item 1's current framing**: the headline currently states
"environmental conditioning is a real trade-off... it helps most where the shared curve already
fits worst, and substantially degrades accuracy where it already fits well" -- framing BOTH halves
(the Q4 help and the Q1 harm) as effects of environmental conditioning specifically. This check
shows that framing is only half right. Q4's benefit is genuinely, if partially, attributable to
real environmental signal (confirmed directly, not just consistent with a plausible account). Q1's
harm is NOT attributable to environmental conditioning specifically -- it happens at close to the
same magnitude when the "environmental" values fed to the model are pure noise, which means it is
better described as an inherent cost of giving any flexible model enough capacity to have output
variance on rows the baseline already fits well, not a specific cost of using real environmental
data. The "trade-off" is real, but it is a trade-off from adding model flexibility in general, not
specifically from environmental conditioning -- a materially different claim from what the current
draft states.

**What this does not establish**: this is one Set (Set3), one model family (XGBoost, fixed
config), pooled across 5 folds at a single permutation seed per fold, not a full permutation-null
distribution (e.g. not repeated across many random seeds to build a null band the way
`compute_residual_morans_i`'s 999-permutation Moran's I test does elsewhere in this project). The
DNN/PINN versions of RQ2a were not re-tested this way. The direction and rough magnitude of the
Q1-vs-Q4 asymmetry is clear and consistent across both cohorts, which is reassuring, but this
should be read as a single, well-targeted control confirming the mechanism, not an exhaustive
robustness sweep.
